\documentclass{article}
\usepackage[T1]{fontenc}
\usepackage{lmodern}
\usepackage{graphicx}
\usepackage{amsmath,amssymb}
\usepackage[a4paper,margin=2cm]{geometry}
\usepackage[absolute,overlay]{textpos}
\setlength{\TPHorizModule}{1mm}
\setlength{\TPVertModule}{1mm}
\usepackage[indonesian]{babel}
\usepackage{hyperref}
\setlength{\emergencystretch}{2em}
% unit: Halaman 17

\begin{document}

% === Halaman 17 (python/native) ===

17

yang dalam hal ini nilai Pw dan Qw (dalam heksadesimal) berbeda-beda bergantung pada w 

sebagai berikut [STA98]: 


w 

16 

32 

64 

Pw 

B7E1 

B7E15163 

B7E151628AED2A6B 

Qw 

9E37 

9E3779B9 

9E3779B97F4A7C15 



Konstanta Pw dan Qw didasarkan pada representasi bilangan alam e dan  dalam biner,  

Pw = Odd[(e – 2)2w]  

Qw = Odd[( – 1)2w]  


yang dalam hal ini, 

e = 2.718281828459…  

 = 1.618033988749… = 

2

5

1+



Odd adalah fungsi yang memberikan integer ganjil terkecil yang 

dekat dengan input yang diberikan

\end{document}
